
\section{Dataset construction}
\label{sec:data}

\subsection{A twelve-market corpus of listings}

The corpus is built from sold residential listings scraped from a major online
real estate platform across twelve metropolitan markets that span the price and
geographic range of the United States housing stock. The markets run from San
Francisco and New York at the expensive coastal end, through Boston, Seattle,
Washington, Denver, and Portland in the middle, to Chicago, Atlanta, Phoenix,
Dallas, and Philadelphia toward the more affordable interior and south. Each
listing carries the free-text description an agent wrote, the asking price the
platform reports for the property at the time of the crawl, and the structured
attributes a hedonic model expects: the bedroom and bathroom counts, the living
area, the lot size, the year built, and the property type. Each listing also
carries a property-level latitude and longitude rather than a coarse centroid. Two
features of this outcome we are explicit about. The price is
the listing's asking price rather than a realized transaction, and it is set by the
same agent who writes the description and at the same moment. Part of any
text-price association is therefore the agent's own pricing of the language rather
than the market's, a simultaneity that no spatial adjustment removes. It limits the
causal reading of the estimates along an axis orthogonal to the spatial-confounding
question this paper takes up. Because the corpus is a single cross-sectional crawl,
the outcome itself carries no calendar-time variation that pooling could confound.
The property's most recent prior-sale date, which we recover from the listing page
for every property and release with the data, enters only as a covariate. It is in
any case essentially uncorrelated with the asking price within each market: the
absolute correlation is no larger than $0.20$ and the $R^2$ is below $0.04$.
Table~\ref{tab:corpus-12} summarizes the twelve markets with their analysis sizes
and median prices.

\begin{table}[H]
\centering
\caption{The twelve-market corpus. \emph{Listings} is the per-market analysis
sample after deduplication to one row per listing and removal of rows missing a
required field; the markets together contribute $69{,}204$ listings. Median sale
price and median description length are computed on the analysis sample. The
panel ranges over more than an order of magnitude in size and roughly fivefold in
price.}
\label{tab:corpus-12}
\small
\begin{tabular}{lrrr}
\toprule
Market & Listings & Median price & Median words \\
\midrule
New York      & $15{,}255$ & \$975{,}000   & $183$ \\
Dallas        & $8{,}008$  & \$439{,}000   & $161$ \\
Phoenix       & $7{,}087$  & \$514{,}990   & $117$ \\
Philadelphia  & $7{,}030$  & \$255{,}000   & $134$ \\
Atlanta       & $6{,}135$  & \$384{,}000   & $178$ \\
Chicago       & $5{,}577$  & \$374{,}900   & $129$ \\
Denver        & $5{,}280$  & \$525{,}000   & $207$ \\
Portland      & $4{,}337$  & \$559{,}000   & $166$ \\
Washington    & $4{,}015$  & \$599{,}000   & $187$ \\
Seattle       & $2{,}884$  & \$750{,}000   & $152$ \\
Boston        & $2{,}610$  & \$999{,}000   & $137$ \\
San Francisco & $986$      & \$1{,}295{,}000 & $174$ \\
\midrule
Total         & $69{,}204$ & & \\
\bottomrule
\end{tabular}
\end{table}

\subsection{Deduplication, outliers, and a robust winsorizer}

The scraper issued repeated requests against the same listing during collection.
The raw rows therefore resolve to a smaller set of distinct listings once they are
deduplicated on the listing identifier. Every number in the paper is computed
on the one-row-per-listing analysis sample, to keep sibling copies of a listing out
of the training folds of any cross-fitted procedure. We drop transactions outside a
plausible price band to remove data-entry errors and non-arm's-length transfers.
We also guard the structured covariates with a per-column robust winsorizer that
clips each feature at its median plus or minus five median absolute deviations and
imputes the rare non-finite entry at the median. The winsorizer matters more than
its routine description suggests. A single parcel-join mistake, a year-built value
tens of standard deviations from any plausible construction date, was enough to
make the partialling-out estimator in one market ill-conditioned and to drive its
coefficient to an impossible magnitude with a deceptively small standard error.
Clipping the true outliers without touching the well-behaved rows restored that
market to the range of the others, while leaving every well-behaved market
unchanged. The clip is applied uniformly across the twelve markets and is the only
manipulation of the structured covariates we perform.

\subsection{Confounders}

The conditioning set is assembled to make the backdoor adjustment of
Section~\ref{sec:methods} credible at the resolution where language and price are
both determined. Alongside the structured property attributes, we attach the
property's location through its coordinates and its ZIP code. We join
neighborhood context by spatial position rather than by a coarse identifier. Census
demographics come from the American Community Survey five-year estimates at the
block-group level, the finest geography the survey supports. They contribute median
household income, educational attainment, racial composition, median home value,
and median gross rent for the block group containing each listing. Recent crime is
counted within a five-hundred-meter radius of each property from municipal incident
records, separated into violent, property, and quality-of-life categories through a
cross-city offense crosswalk. Local amenities are counted in the same radius from
an open geographic database, covering food and retail, services, recreation,
transit, and education. The assembled set runs to roughly three dozen controls per
market. The spatial block within it, the coordinates together with the ZIP
indicators, is the group the sensitivity analysis of Section~\ref{sec:sensitivity-12}
singles out, because it is the one most directly implicated in the
text-as-spatial-proxy concern. Appendix~\ref{app:data_collection} documents the
scrape, the geocoding, and the census, crime, and amenity joins.

\subsection{Text preprocessing and the two encodings}

Before encoding, each description is stripped of markup and boilerplate,
lowercased, and cleared of explicit street addresses. The addresses are replaced
with anonymization tokens, so that a model cannot recover location by reading an
address out of the text rather than out of the language. The preprocessed text
then feeds the two encodings that Section~\ref{sec:methods-channels} defines. The
Doc2Vec uniqueness score that drives the Shen-Ross channel measures how far a
listing's language sits from its market's corpus norm. The sentence-transformer
embedding that drives the Baur channel is the $768$-dimensional representation of
\texttt{all-mpnet-base-v2} \citep{reimers2019sentence}, from which we take the
standardized leading principal component as the text treatment. Robustness to the
choice of encoder is taken up directly in Section~\ref{sec:cross-method-12}, where
a Doc2Vec uniqueness axis and the sentence-BERT component recover a concordant
per-market signal, so the effect does not hinge on one encoder's geometry. The
deconfounding and erasure diagnostics of Sections~\ref{sec:leace-12}
and~\ref{sec:decomposition} operate on the same mpnet embedding, so that what they
remove is the geography the Baur treatment itself carries.
